% Part 4: course-ordered continuation, 2026-10-14 to 2026-10-16.
\section{从隐藏节点时序到设备间直接通信}
\label{ch:d2d}

\subsection{隐藏节点条件下的退避推演}

分布式协调功能（Distributed Coordination Function, DCF）的计算题不能只画一条代表全网的“信道忙闲线”。每个站点掌握的信道状态取决于它能够听到哪些发送端。同一时刻，接入点可能正在接收，某个隐藏站点却仍把信道判断为空闲。因此，时序分析必须先确定拓扑，再分别记录各节点的物理侦听、网络分配向量（Network Allocation Vector, NAV）、等待间隔和退避计数。

\subsubsection{不同包长碰撞后的不对称状态}

设两个站点同时发送，但A的数据帧较短，B的数据帧较长。A发送期间不能同时接收；A结束发送后，可能只听到B帧的尾部。虽然检测到了无线能量，A并没有从头接收到可正确解码的帧，不能根据这个尾部恢复可靠的Duration字段。因而，在错误接收所对应的情形下，A待信道重新空闲后使用扩展帧间间隔（Extended Interframe Space, EIFS）。B则需要根据确认超时判断本次发送失败，再进入重传竞争。

这里有两条不同的因果链：A的额外等待来自接收错误或无法正确解码的活动，B的重传来自未收到确认帧（Acknowledgment, ACK）。不能因为两者参加了同一次碰撞，就让两者在完全相同的时刻进入同一种等待状态。若C在碰撞期间收到上层待发分组，还须根据C自身是否成功读取首部，决定它使用NAV还是错误接收后的保护等待。待C先取得发送机会后，能够正确接收C首部的其他站点设置相应NAV，并冻结尚未完成的退避。

\subsubsection{两包连续发送与后退避}

考虑图\ref{fig:p4-hidden}的拓扑。A与B相互隐藏，C能够听到A、B和接入点，接入点也能与所有站点通信。B先收到两个待发分组，第一帧发送成功后，不能直接把第二帧接在后面独占信道，而要执行后退避（Post-backoff）。如果只研究每站一帧的例子，成功发送后的后退避可能不影响所关心的结果；当同一站点还有第二帧时，它就成为决定发送顺序的关键状态。

\begin{figure}[htbp]
\centering
\begin{tikzpicture}[scale=0.95,every node/.style={font=\small}]
\node[draw,circle,minimum size=9mm] (a) at (0,0) {A};
\node[draw,circle,minimum size=9mm] (c) at (3,0) {C};
\node[draw,circle,minimum size=9mm] (b) at (6,0) {B};
\node[draw,rounded corners,minimum width=16mm,minimum height=9mm] (ap) at (3,2) {AP};
\draw (a)--(c)--(b); \draw (a)--(ap)--(b); \draw (c)--(ap);
\draw[dashed] (0,-1)--(6,-1);
\node[fill=white] at (3,-1) {A与B之间无可侦听链路};
\end{tikzpicture}
\caption{隐藏节点例子的原创拓扑示意。实线表示题设中的相互可听关系，不表示节点必定同时发送。}
\label{fig:p4-hidden}
\end{figure}

推演时应先标出上层分组到达、已经给定的发送开始时刻和帧长，再推导其余动作。尤其要区分“收到待发分组”和“已经开始发送”：前者还可能需要等待DCF规定的空闲间隔，后者则是时间轴上的既定事件。接下来分别画出ACK、NAV与退避，最后才比较谁的计数器先归零。

第三个例子以下列已经确定的中间计数状态为起点，展开剩余退避约束，并不重新给定完整的绝对到达时间题干：B在A的数据发送及ACK附近的危险边界前，共经历13个可递减时隙。与此同时，C原来抽取7，在A接入前已递减2，因此在A与AP完成交换后还剩5。令B初始后退避整数为$b$，则在这个重新竞争的阶段，其剩余值为
\begin{equation}
 r_B=b-13,\qquad r_C=7-2=5.
\end{equation}
B要在C之前发送，要求$r_B<5$；B又不能在AP准备发送ACK或仍不能接收的边界上归零，要求$r_B\geq1$。两条约束合并为
\begin{equation}
 1\leq b-13\leq4
 \quad\Longrightarrow\quad b\in\{14,15,16,17\}.
\end{equation}
这不是把ACK持续时间也折算成B必须额外抽取的时隙。B能够听到AP，一旦ACK开始，B的物理侦听就会使剩余退避冻结；ACK结束并经过分布式帧间间隔（DCF Interframe Space, DIFS）后，才继续消耗剩余时隙。

取$b=13$的问题在于计数器过早归零：A的数据发送虽然可能刚结束，AP尚处于准备回复ACK或发送ACK的阶段。半双工（Half-duplex）接入点不能一边发ACK、一边接收B的新数据。取$b=18$则留下5个时隙，使B与C同时归零、同时发送。更大的值可能让C先发送而B随后成功，但不能满足“B的两帧都先于C完成”的限定。

\subsubsection{只限制碰撞次数时，可行区间如何改变}

第四个例子保留同一隐藏节点拓扑，但先给定B开始发送，A稍后开始发送，并使两帧在同一时刻结束。C在重叠发生之前已经读取B的首部，因此即使后续数据碰撞，C仍可按先前读出的Duration维持预期交换的NAV。A、B均因没有收到ACK而进入失败后的竞争。

继续使用例题给出的中间计数状态：A的给定退避整数为1，C的为5；B看不到A的发送，故仍会递减，并在AP的ACK使其冻结前已经消耗16个完整可递减时隙。这里把16作为已给出的中间结果，随后完整求解可行退避集合。于是
\begin{equation}
 r_B=b-16,\qquad r_C=5.
\end{equation}
若B先发送，必须有$1\leq r_B\leq4$，得到$b=17,18,19,20$。若B、C同时归零，则$r_B=5$，对应$b=21$，应当排除。若允许C先发送，$r_B\geq6$也可行；C发送时，B听到C并冻结，C完成交换之后，B继续剩余退避。

因此，在例题限定的抽样范围$0\leq b\leq31$内，可行集合为
\begin{equation}
 \mathcal{B}=\{17,18,19,20\}\cup\{22,23,\ldots,31\}.
\end{equation}
$b=16$不能取，原因仍是AP的收发切换与ACK保护，而非A的数据帧本身尚未结束。两个例子的共同方法是“先保护不可侦听的发送，再比较可共同侦听阶段的剩余计数”；不同答案来自发送顺序约束不同。这里的31是例题采用的候选上限，不应据此取消一般DCF在失败后调整竞争窗口的规则。

\subsection{设备到设备网络的动机与工作方式}

设备到设备通信（Device-to-Device Communication, D2D）允许邻近设备通过直接无线链路交换数据。自组织网络（Ad hoc Network）进一步强调：节点可以按需要组成网络，不必先部署一套固定基础设施。灾害救援中既有网络可能损坏；临时活动中可能需要快速建立局部通信；个人手机、电脑、耳机和物联网标签也常需要近距离互联。直接通信减少了对固定设施的依赖，部署较快，且能够利用短距离链路降低发射功率。

纯自组织网络（Pure Ad hoc Network）的数据完全在参与设备之间传递。混合网络（Hybrid Network）则同时保留直接通信与基础设施连接，某个节点或网关（Gateway）负责把局部网络接入互联网。网关的存在不意味着所有局部数据都必须绕行互联网；同一网络可以让近邻直接交换，同时让跨网络业务经过网关。

在蜂窝网络辅助的D2D场景中，基站仍可能管理资源或提供外部连接，但局部内容分发可由设备协作完成。例如，许多邻近用户需要同一内容时，基站向少量设备传送，随后由这些设备向附近用户转发，就可能降低基站重复发送的负担。这种收益以设备距离近、链路可靠、资源协调适当以及设备愿意转发为前提。

\subsubsection{空间复用的含义}

空间复用（Spatial Reuse）是指在彼此干扰足够小的不同空间区域重复使用同一频率资源。设一个服务区域可使用$K$个信道，每个信道在给定业务负载下支持约$u$个用户，则单个覆盖区域的简单容量估计为$Ku$。把服务区域划分为多个低功率区域，并在满足干扰约束的位置重复使用同一组信道，有机会增加同时服务的总用户数。实际容量不会无条件按区域数线性增长，因为边缘干扰、调度和复用间距都会限制增益。

D2D把这一思想带到设备附近。较短的传输距离通常允许较低的发射功率，较低功率使强干扰的范围缩小，于是远处的另一对设备可能同时使用同一信道。若设备并不邻近，必须提高功率才能完成直接通信，空间复用收益便可能减弱。

\begin{figure}[htbp]
\centering
\begin{tikzpicture}[every node/.style={font=\small}]
\draw[dashed] (0,0) circle (1.25); \draw[dashed] (4.6,0) circle (1.25);
\node[draw,circle] (u1) at (-0.45,0) {$S_1$};
\node[draw,circle] (v1) at (0.45,0) {$D_1$};
\node[draw,circle] (u2) at (4.15,0) {$S_2$};
\node[draw,circle] (v2) at (5.05,0) {$D_2$};
\draw[->,thick] (u1)--(v1); \draw[->,thick] (u2)--(v2);
\node at (0,-1.7) {近距离、低功率链路};
\node at (4.6,-1.7) {可复用同一信道};
\node at (2.3,1.7) {是否能并发取决于接收端的实际干扰};
\end{tikzpicture}
\caption{空间复用的原创概念图。虚线仅表示局部干扰范围的简化示意，不是固定的无线传播边界。}
\label{fig:p4-reuse}
\end{figure}

\subsection{业务、设备与移动性的差异}

D2D中的节点并不一定同质。手机和笔记本通常比微型传感器具有更多计算资源、更大的存储空间和不同的电池容量。设备的逻辑责任也可能不同，例如某种技术安排一个设备协调通信，另一些设备只完成简单数据交换。能力较强并不自动意味着承担所有管理工作，角色仍取决于协议和网络组织。

业务目的决定了交付方式。单播（Unicast）面向一个确定目的端；广播（Broadcast）面向广播范围内的所有节点；多播（Multicast）面向指定接收组。地理广播（Geocast）按地理区域选择接收者，例如让某一道路范围内的车辆接收交通告警；任播（Anycast）则只需把数据交给某个合适的组成员，例如若干可提供同一服务的节点中的一个。任播不要求所有成员都收到。

业务要求至少应同时描述比特率、时延和可靠性。以成功交付比例描述可靠性时，可以写成
\begin{equation}
 \rho=\frac{N_{\mathrm{delivered}}}{N_{\mathrm{generated}}}.
\end{equation}
分母是观察期间产生的目标分组数，分子是按所选成功判据交付的分组数；如果应用另有截止时间，还应明确是否只把及时到达的分组算作成功。低速率温度测量与实时告警的要求不同，不能只凭平均流量判断设计优先级。

移动性（Mobility）既包括移动速度，也包括移动模式。游牧式使用往往是位置改变后停留；车辆沿道路高速移动；家庭设备、机器人或随机移动的人群又具有不同规律。速度越高，邻居关系的有效时间可能越短，网络发现和路由维护的工作也可能越频繁。固定节点同样可能因障碍物、故障或休眠而失去可用链路。

\subsection{拓扑、接入和多跳转发}

\subsubsection{用随时间变化的图表示网络}

网络拓扑（Network Topology）可抽象成图
\begin{equation}
 G(t)=\bigl(V(t),E(t)\bigr),
\end{equation}
其中$V(t)$是当前可参与通信的节点集合，$E(t)$是当前可用链路集合。邻居（Neighbor）指能够直接通信的节点。简单模型常把通信范围内的链路视为可靠，把范围外的链路视为不存在；真实系统中，衰落、干扰和接收门限会使边的可用性更加复杂。

通信范围也不能与干扰范围混为一谈。一个信号可能已经弱到无法正确解调，却仍足以破坏另一个接收机正在接收的数据。因而，“不是邻居”不等于“互不干扰”。半双工限制进一步要求节点分时完成发送和接收。

平面拓扑（Flat Topology）中没有强制的层级管理结构，网状拓扑（Mesh Topology）可提供多条替代路径。对$N$个节点的简单无向图，每对不同节点至多对应一条边，因此完全连接时的最大边数为
\begin{equation}
 |E|_{\max}=\frac{N(N-1)}{2}.
\end{equation}
分子$N(N-1)$把每一对节点计算了两次，除以2得到无向边数。一般的平面网并不必然达到这个上限。

树形拓扑（Tree Topology）在连通且无环时恰有$N-1$条边。两节点之间只有一条简单路径，所以路径中的一条边失效即可使某些节点之间失去连通。网状结构若存在冗余路径，单边失效后可能绕行；但冗余连接也意味着更多潜在干扰和状态维护工作。

分簇拓扑（Cluster-based Topology）把节点组织为若干簇，由簇头（Cluster Head）承担簇内协调或数据转发。桥接节点（Bridge Node）连接相邻簇，使数据能够跨簇逐跳传播。分簇可以减小局部管理范围，却也会使簇头和桥接节点成为重要的负载集中点。

\subsubsection{接入负责一跳，路由负责路径}

信道接入（Channel Access）决定某一跳的多个竞争者何时使用共享无线信道；路由（Routing）决定数据经过哪些中继（Relay）到达目的端。即使每一跳都遵守MAC规则，网络仍需要选出可到达目的端的路径。反过来，即使路径已经选好，每一跳也仍需获得接入机会，路径存在并不意味着能立即发送。

网络层（Network Layer）执行路由功能。源节点选择路径或下一跳，中继收到数据后继续选择下一跳；MAC层完成本跳发送规则，物理层（Physical Layer）执行实际无线信号传输。应用层产生的信息向下交付，成功接收后再由目的端协议栈向上交付。不能把中继仅看成“把听到的信号原样放大”：这里讨论的是节点接收分组后进行下一跳转发。

\begin{figure}[htbp]
\centering
\begin{tikzpicture}[every node/.style={font=\small}]
\node[draw,circle] (a) at (0,0) {A}; \node[draw,circle] (c) at (6,0) {C};
\node[draw,circle] (b) at (3,1.6) {B};
\node[draw,circle] (d) at (2,-1.6) {D}; \node[draw,circle] (e) at (4,-1.6) {E};
\draw[->] (a)--node[above]{一跳直达}(c);
\draw[->,thick] (a)--(b); \draw[->,thick] (b)--node[above right]{两跳}(c);
\draw[->,dashed] (a)--(d); \draw[->,dashed] (d)--(e); \draw[->,dashed] (e)--(c);
\node at (3,-2.3) {三跳候选路径：A--D--E--C};
\end{tikzpicture}
\caption{同一源、目的端的不同候选路径。原创图用于说明路径选择，并不规定某条路径总是最优。}
\label{fig:p4-routes}
\end{figure}

多跳通信（Multi-hop Communication）可把一个较远的传输分解为数个近距离传输。源节点因只需到达邻近中继，可能降低自身发射功率；较小的干扰范围又可能改善空间复用。但每增加一跳，就增加一次信道使用、排队与接入过程，也增加中继的接收和发送工作。因此，源节点节能不等于全网总能耗必然降低，跳数较少也不保证路径在所有指标上最优。

链路失效时，若还存在候选路径，路由可转向备用路径。这样的稳健性以节点保持参与和愿意中继为前提。民用设备的拥有者可能不愿为他人承担电池和资源开销，需要激励机制；任务型网络可能由共同任务规定合作责任。后续分析采用中继愿意合作转发的模型。

\subsection{能量约束与分布式设计}

移动设备通常依赖电池；固定在偏远、恶劣环境中的传感器也可能无法接入电源。电池容量受到体积、重量和维护条件限制。若用能量预算$E_{\mathrm{bat}}$和长期平均功率$\bar P$表示，理想化寿命估计为
\begin{equation}
 T_{\mathrm{life}}\approx\frac{E_{\mathrm{bat}}}{\bar P}.
\end{equation}
这一关系说明，降低长期平均功率与增加电池能量都能延长寿命；但在体积已受限制时，协议设计往往只能从减少无效活动入手。

\subsubsection{功率控制与连通性的平衡}

功率控制（Power Control）通过调整发射功率，在可达性、干扰和能耗之间取舍。采用简化路径损耗模型时，接收功率可写为
\begin{equation}
 P_r=K\frac{P_t}{d^{\alpha}},
\end{equation}
$P_t$是发射功率，$d$是收发距离，$\alpha$是路径损耗指数，$K$汇总模型中的其他固定因素。若需要维持相同的接收门限$P_{\min}$，则
\begin{equation}
 P_t\geq\frac{P_{\min}}{K}d^{\alpha}.
\end{equation}
因此，把距离缩短为原来的某个比例，会以该比例的$\alpha$次方改变所需发射功率。这解释了近邻转发的节能潜力，也说明降低功率不能脱离接收门限随意进行。

高功率通常增加可用邻居和冗余边，但会扩大干扰；低功率减小干扰，却可能使网络变得稀疏甚至分裂。合理目标是以足够低的功率保持任务需要的连通性和可靠性，而非把所有节点始终设为最大功率或最小功率。

\subsubsection{活动调度与休眠深度}

活动调度（Activity Scheduling）决定何时保持设备清醒、何时进入休眠（Sleep）。无线接口即使没有发送，也可能因持续侦听而消耗显著能量。关闭不必要的收发和计算组件可以节能，但睡着的节点不能立即接收数据，还可能暂时从可用拓扑中消失。

浅休眠保留较多状态，唤醒较快，但功耗相对高；深休眠关闭更多组件，功耗更低，但恢复工作通常需要更长时间。选择休眠深度时，应同时考虑下次预期活动时刻、唤醒代价、允许的接收时延和网络对该节点的转发依赖。频繁进入、退出深休眠未必比保持短暂浅休眠更省能。

拓扑、信道接入、多跳路由和能量消耗构成D2D设计中互相影响的四个问题。自治自组织网络通常需要分布式算法（Distributed Algorithm），让节点根据本地观测与有限交换作决定；集中式算法（Centralized Algorithm）则把信息和决策集中于管理者。分布式并不等于没有任何协调，而是不能把所有正确运行条件都建立在一个永远可达的全局控制者之上。自组织Wi-Fi、蜂窝辅助D2D、车联网、蓝牙和体域网（Body Area Network, BAN）体现了不同的角色、移动性和能量约束；传感器网络则把低功耗要求推到了设计中心。

\section{无线传感器网络与低功耗介质访问}
\label{ch:wsn}

\subsection{从物理测量到网络数据}

无线传感器网络（Wireless Sensor Network, WSN）由分布于环境中的传感节点组成，是物联网（Internet of Things, IoT）的重要组成部分。节点测量温度、压力、气体浓度、位置或其他物理量，并把有用信息交给汇聚节点（Sink）。汇聚节点连接局部传感网络与外部系统，承担数据出口，有时也负责向节点下发控制信息。

大范围森林火灾监测说明了其典型工作链：各位置的传感器定期测温，本地处理识别异常，必要时报告告警；远端系统收到告警后再组织响应。偏远或寒冷地区的环境监测同样依赖长期电池供电，维护困难使通信效率和寿命尤为重要。

\subsubsection{传感节点的基本组成}

传感板（Sensing Board）把物理量转换成可处理的测量值，信号处理单元（Signal Processing Unit）进行分析、筛选或事件检测，无线接口（Wireless Interface）负责数据交换。节点还可以控制执行器（Actuator），把计算结果转化成物理动作。例如，超过门限的温度测量既可以触发远端告警，也可以在本地触发某种控制动作。

\begin{figure}[htbp]
\centering
\begin{tikzpicture}[every node/.style={font=\small},box/.style={draw,rounded corners,text width=23mm,align=center,minimum height=11mm}]
\node[box] (sense) at (0,0) {传感板\\物理量采样};
\node[box] (proc) at (3.2,0) {信号处理\\事件检测};
\node[box] (radio) at (6.4,0) {无线接口\\测量与告警};
\node[box] (act) at (3.2,-1.8) {执行器\\本地控制};
\draw[->] (sense)--(proc); \draw[<->] (proc)--(radio); \draw[->] (proc)--(act);
\node at (0,-1.8) {电源支持各模块};
\end{tikzpicture}
\caption{传感节点的信息处理链原创示意。模块可以集成于同一设备，执行器并非每个节点都必须具备。}
\label{fig:p4-sensor}
\end{figure}

物理现象通常随时间和空间连续变化，而网络只能处理有限测量。设传感器位于$\boldsymbol{x}_i$，在时刻$t_k$采样，则样本可记为
\begin{equation}
 y_i[k]=s(\boldsymbol{x}_i,t_k),
\end{equation}
其中$s$是所观察的物理场。时间采样间隔决定变化捕捉的细致程度，空间部署决定不同位置的覆盖程度。本地处理与外部系统可利用这些离散数据重建或估计物理变化；仅有少量采样点并不自动保证重建精确，采样密度应适应现象变化的尺度。

\subsubsection{传输介质与低速率特征}

传感器的无线接口可以采用射频（Radio Frequency, RF）、红外或光通信。红外、光通信对视距（Line of Sight, LOS）和方向通常更敏感，障碍物可能直接遮断链路。射频常更适合存在遮挡的环境，但仍受传播条件影响。微型“智能尘埃”（Smart Dust）体现了把感知、处理和通信集成到微小设备中的思路，不能据此假定所有传感器都使用同一种传输介质。

本课程采用的典型低速率节点尺度包括约5--150米的射频覆盖示例、约1--500 kbit/s的数据速率量级，以及约1 mW的低发射功率示例。这些数值用来说明低功耗短报文设计，而非所有产品或所有IEEE 802.15.4物理层的统一性能。以分贝毫瓦（dBm）表示功率时，
\begin{equation}
 P_{\mathrm{dBm}}=10\log_{10}\left(\frac{P}{1\,\mathrm{mW}}\right),
\end{equation}
所以1 mW对应0 dBm。小型电池、低成本和小尺寸共同限制节点资源，协议也因此不能照搬高吞吐量终端的设计优先级。

\subsection{汇聚业务、寻址与协议栈}

传感节点可以上传原始测量，也可以先在本地处理，只上传压缩后的摘要或检测到的事件。数据常由许多节点流向少数汇聚节点，构成以汇聚为中心的上行（Uplink）业务；配置、查询、时间信息或控制指令则沿下行（Downlink）方向传播。上下行业务量通常不对称，但下行流量较少并不意味着可以忽略，因为它可能决定采样、同步与控制能否正确进行。

当大量节点传输短载荷时，寻址开销十分明显。若应用数据只有$L$字节，地址及其他协议开销共$H$字节，则报文有效比例为
\begin{equation}
 \eta=\frac{L}{L+H}.
\end{equation}
当$L$很小时，固定首部$H$占比就很大。短地址能减轻局部传输负担，但仍须保证所使用范围内的唯一性，并处理与外部寻址体系的衔接。不能只从地址更短就推出整个系统不需要地址管理。

时间同步（Time Synchronization）让不同节点的测量可在时间上比较，也支持协同休眠和时隙调度。自定位（Self-localization）则给测量加上空间含义：温度值相同，但出现位置不同，可能意味着完全不同的事件。感知数据是否有用，不仅取决于它是否到达，还取决于“何时、何地测得”是否可解释。

\subsubsection{应用如何形成设计要求}

环境监测、滑坡预警、动物跟踪、建筑管理、车辆系统、水质和土壤监测，都利用了分散节点在现场收集信息的能力。智能电网（Smart Grid）需要观察发电、负载和设备状态，以支持调节与控制；可再生能源的变化使及时、分布式测量具有价值。土壤污染、水中微生物和农业条件的监测又强调长期部署、覆盖与维护成本。相同的感知网络框架可以服务不同应用，但采样周期、时延容忍度和可靠性目标应随应用改变。

商用传感平台体现了模块化实现：同一平台可配置不同传感板，选择不同无线接口，并提供多种休眠模式。气体、温度、液体、重量和压力等测量能力来自不同传感器模块。接口种类、发射功率和覆盖距离应逐项理解，不应把一个平台不同配置的最佳指标同时当成单一节点的固定能力。

\subsubsection{协议层之间的配合}

物理层负责信号传输；MAC层负责接入共享介质和组织收发活动；网络层负责路由，并适应因移动、故障或休眠而变化的拓扑。较高层还可能提供缓存（Caching）、查询（Query）和数据聚合（Data Aggregation），把多个测量整合成应用需要的信息，再与外部数据库或业务系统连接。

休眠会使“无线链路在地理上可达”和“节点此刻可接收”成为两件事。路由若选择了正在深睡的中继，MAC即使找到空闲信道，也无法立即完成交付。同步、定位和能量管理会跨越多个层次发生影响，分层用于明确职责，而不能消除这些协同需求。

\subsection{MAC设计的主要目标与能量浪费}

无线传感器网络的MAC尤其重视能量效率（Energy Efficiency）、可扩展性（Scalability）和拓扑适应性（Adaptability）。节点增多后，协议不应因控制负担过大而失效；节点休眠或重新加入时，接入机制也应继续运行。信道利用率、时延、吞吐量（Throughput）和公平性（Fairness）同样重要，只是其优先级取决于应用。温度低频采样与视频传输的吞吐量需求不同；突发告警即使平均速率很低，也可能要求很短的交付时延。

分布式接入意味着节点根据自己的观察决定行动。多个节点可能同时发现信道空闲并同时发送，因此“独立作出合理判断”并不等于“全网必然无碰撞”。MAC需要通过竞争分散或预先调度来减少这类同时行动。

\subsubsection{功耗状态与活动时间}

无线接口可处于发送、接收、空闲侦听和休眠状态。低功耗设计中的重要事实是：空闲侦听的功率可能与接收功率处于相近量级，而远高于休眠功率。总能量可按各状态停留时间估计为
\begin{equation}
 E\approx P_{\mathrm{tx}}t_{\mathrm{tx}}+
 P_{\mathrm{rx}}t_{\mathrm{rx}}+
 P_{\mathrm{idle}}t_{\mathrm{idle}}+
 P_{\mathrm{sleep}}t_{\mathrm{sleep}}.
\end{equation}
式中每项都是“该状态的功率乘以停留时间”。如果绝大多数时间都没有有效业务，却保持空闲侦听，那么即使每次发送已经很短，长期电量仍可能被侦听耗尽。CPU、传感板以及模式切换的消耗也应在完整预算中计入，上式仅把无线状态的主线展开。

四类典型浪费分别是：空闲侦听（Idle Listening），即没有有用接收却持续开机；碰撞（Collision），即已花费收发能量但数据不能正确交付；协议开销（Overhead），即控制报文、首部和维护活动消耗的资源；旁听（Overhearing），即接收并处理本不发给自己的报文。节能MAC不能只减少碰撞，还应减少无效清醒和不必要接收。

\subsection{基于调度的接入：LEACH}

基于调度的MAC（Schedule-based MAC）预先安排谁在何时发送。

基于竞争的MAC（Contention-based MAC）则由节点在业务出现时竞争接入。调度可减少碰撞并让节点知道何时休眠，但需要维护日程；竞争更灵活，却可能产生等待、碰撞和重试。

低能耗自适应分簇层次协议（Low-Energy Adaptive Clustering Hierarchy, LEACH）展示了分簇与调度结合的思路。普通节点向本簇簇头报告数据，簇头承担更多收集和转发工作。簇内采用时分多址（Time Division Multiple Access, TDMA），把重复的时间帧划成时隙，每个节点在自己的时隙发送，其他时间则有机会休眠。

\begin{figure}[htbp]
\centering
\begin{tikzpicture}[every node/.style={font=\small},x=1.2cm,y=1cm]
\foreach \x/\s in {0/S_1,1/S_3,2/S_5,3/S_1,4/S_3,5/S_5}{
\draw (\x,0) rectangle (\x+1,0.8); \node at (\x+0.5,0.4) {$\s$};}
\draw[<->] (0,1.15)--node[above]{第一个时间帧}(3,1.15);
\draw[<->] (3,1.15)--node[above]{下一个时间帧}(6,1.15);
\node[align=center] at (3,-0.65) {簇头按表接收；普通节点主要在自己的时隙活动};
\end{tikzpicture}
\caption{LEACH簇内TDMA的原创时隙示意。时隙标签仅用来说明重复调度。}
\label{fig:p4-leach}
\end{figure}

在理想同步与正确调度下，同一簇内不同节点不会占用同一个发送时隙，因此可避免簇内冲突。不同簇可能同时工作，课程采用码分多址（Code Division Multiple Access, CDMA）说明簇间区分：不同码使接收端能够区分目标信号和其他同时存在的信号。其直观意义是选择所关注的信号，而不是让其他信号在物理上消失；可分离性仍取决于码、功率和干扰条件。

簇头需要接收多个节点的数据并向外转发，能耗通常高于普通节点。轮换簇头（Cluster-head Rotation）把较重的职责在节点之间轮流分配，有助于避免固定簇头过早耗尽电池。轮换并不消除协调成本，但能改善负担分布。

静态时隙安排面临可扩展性问题。若时间帧已有三个固定时隙，新节点加入却没有空余时隙，就必须重新调度、延长帧，或降低某些节点的发送频率。设每个时隙长$T_s$，一个簇安排$N$个发送节点且每帧各发送一次，忽略其他开销时，
\begin{equation}
 T_{\mathrm{frame}}=NT_s,\qquad
 f_{\mathrm{report}}=\frac{1}{NT_s}.
\end{equation}
保持时隙时长不变而增加节点，会增加每节点的报告间隔。若坚持固定帧长，就需要其他资源调整，不能凭空增加无碰撞发送机会。

簇数量增加也可能使可用码资源不足；重复使用相同或相互干扰较大的码，需要考虑空间分隔与簇间干扰。LEACH由此说明了调度类方案的基本权衡：低冲突、可休眠的日程，与节点变化、调度维护和资源扩展之间存在张力。

\subsection{IEEE 802.15.4的分层位置、设备与拓扑}

IEEE 802.15.4定义低速率无线个人区域网（Low-Rate Wireless Personal Area Network, LR-WPAN）的物理层和MAC层。Zigbee基于这些底层能力提供更高层的网络与应用组织。两者并不是同一个名称的两种写法：物理层和MAC层可以支撑不同上层协议，上层的路由能力也不应全部归到IEEE 802.15.4本身。

全功能设备（Full-Function Device, FFD）具有较完整的协议功能，可承担协调器等角色；精简功能设备（Reduced-Function Device, RFD）面向较简单的终端任务，降低实现复杂度，通常通过FFD交换数据。具体多跳路由由所采用的上层网络协议完成，不能把“具有全功能MAC能力”直接等同于“IEEE MAC已经规定完整路由算法”。

个人区域网协调器（PAN Coordinator）管理网络成员，记录节点的加入与离开，并可提供同步参考。星形拓扑（Star Topology）以协调器为中心，外围终端可采用较简单的设备。对等拓扑（Peer-to-Peer Topology）允许有能力的节点之间直接通信，扩展连接结构；簇树拓扑（Cluster-tree Topology）则把多个局部星形或簇连接起来，由全功能节点承担中间协调角色，使较大覆盖范围与简单叶节点并存。

这些拓扑影响能量负担与故障后果。星形中心易于集中管理，但外围节点需要能够到达中心；更复杂的连接结构允许逐跳扩展，却增加转发与维护工作。后续MAC分析以传感节点和汇聚或协调节点之间的交换为基本场景，并特别关注节点如何安排清醒时间。

\section{IEEE 802.15.4的信标与竞争接入过程}
\label{ch:154}

\subsection{物理层频带与短分组}

本课程讨论的经典低速率物理层包括868/915 MHz附近的频段与2.4 GHz频段。频带中信道的可用带宽影响能够承载的符号速率，调制和编码又决定有效比特率，因此带宽是速率约束的重要组成部分，但不能单凭频率更高就认定速率必然更高。课件把较窄信道与较低速率、较宽信道与较高速率并列，用于说明这一设计关系。

短报文适合少量测量与控制信息。经典帧结构中，物理层前导、定界与长度等字段合计6字节，随后承载最长127字节的物理层服务数据单元（PHY Service Data Unit, PSDU）。这127字节承载的是MAC帧，其中还包含MAC首部和校验等开销，所以应用有效数据小于这个上限。讨论“载荷”时必须说明所属层次，物理层载荷与应用载荷不能直接互换。

\subsection{信标使能模式与超帧}

信标使能模式（Beacon-enabled Mode）中，PAN协调器周期性发送信标（Beacon），为其他节点提供时间参考。由信标开始组织的超帧（Superframe）把活动区间划分成不同接入阶段。先出现竞争接入期（Contention Access Period, CAP），节点采用时隙化竞争；随后可安排无竞争期（Contention-Free Period, CFP），为特定通信提供分配的时隙。

\begin{figure}[htbp]
\centering
\begin{tikzpicture}[every node/.style={font=\small},x=1cm,y=1cm]
\draw (0,0) rectangle (0.7,1); \node[rotate=90] at (0.35,0.5) {信标};
\draw (0.7,0) rectangle (4.5,1); \node[align=center] at (2.6,0.5) {CAP\\时隙CSMA/CA};
\draw (4.5,0) rectangle (7,1); \node[align=center] at (5.75,0.5) {CFP\\分配时隙};
\draw[dashed] (7,0) rectangle (10,1); \node[align=center] at (8.5,0.5) {非活动区间\\可休眠};
\draw (10,0) rectangle (10.7,1); \node[rotate=90] at (10.35,0.5) {信标};
\draw[<->] (0,-0.5)--node[below]{信标间隔}(10,-0.5);
\end{tikzpicture}
\caption{超帧的原创结构示意。图中长度仅表达先后关系；CFP与非活动区间是否存在由配置决定。}
\label{fig:p4-superframe}
\end{figure}

同步意味着节点共享时隙边界，并不意味着CAP由协调器逐个批准发送。CAP内部仍由节点依据算法分布式竞争。CFP则通过分配时隙减少竞争，并让节点知道何时必须清醒。

超帧可留出全网非活动区间，此时包括协调器在内的设备都可进入低功耗状态。单个节点也可以在自己不参与通信的时段休眠。这会把一部分负担转移到缓存和时延：如果目的节点睡着，发送方或协调器必须保存待交付数据，直到合适的唤醒机会。缓存容量有限，等待过久可能导致丢包；因此，休眠节能、存储需求与交付时延需要共同配置。

\subsection{先理解外层：接入许可与停止等待}

时隙载波侦听多路访问与碰撞避免（Slotted Carrier Sense Multiple Access with Collision Avoidance, Slotted CSMA/CA）可以按内外两个逻辑模块理解。外层负责“已经有一帧待发时，发送及确认是否成功”；内层负责“这一次能否取得信道”。先把内层当作只有两种输出的模块，就能把整体流程看清。

外层从初始化发送失败计数开始，然后调用内层。若内层返回允许发送，节点发送数据，并按停止等待（Stop-and-Wait）方式等待ACK。收到正确ACK，本帧成功；若ACK超时，增加失败计数，在尚有尝试预算时重新调用内层，否则报告失败。若内层直接返回接入失败，则本次不能进入数据发送，外层也结束为失败。

\begin{figure}[htbp]
\centering
\begin{tikzpicture}[every node/.style={font=\small},box/.style={draw,rounded corners,align=center,text width=38mm,minimum height=9mm}]
\node[box] (init) at (0,0) {待发帧；初始化发送计数};
\node[box] (inner) at (0,-1.4) {调用内层接入过程};
\node[box] (send) at (0,-2.8) {允许发送：发送并等待ACK};
\node[box] (check) at (0,-4.2) {未获ACK：累计失败次数};
\node[box] (retry) at (0,-5.6) {仍有发送尝试预算？};
\node[box] (fail) at (5,-1.4) {接入失败：结束};
\node[box] (ok) at (5,-2.8) {收到ACK：成功};
\node[box] (drop) at (5,-5.6) {预算耗尽：结束};
\draw[->] (init)--(inner); \draw[->] (inner)--(send); \draw[->] (inner)--(fail);
\draw[->] (send)--(check); \draw[->] (send)--(ok); \draw[->] (check)--(retry);
\draw[->] (retry)--node[above]{否}(drop);
\draw[->] (retry.west)--(-3,-5.6)--(-3,-1.4)--node[above]{是}(inner.west);
\end{tikzpicture}
\caption{按接入与确认两个层次重新绘制的外层流程。接入失败和发送后未获ACK是不同事件。}
\label{fig:p4-outer}
\end{figure}

这一划分防止把所有失败都叫作碰撞。内层可能连续发现信道忙，完全没有发出数据，就用尽了接入尝试；外层也可能已发出数据，但没有获得确认。后者可能来自碰撞或其他接收失败。仅从“没有ACK”通常不能唯一判定物理原因。

\subsection{内层：随机退避、两次CCA与重试}

\subsubsection{三个变量的不同职责}

内层使用退避次数计数$NB$、竞争窗口$CW$和退避指数（Backoff Exponent）$BE$。在此算法中，$CW$表示还需要成功完成多少次空闲信道评估（Clear Channel Assessment, CCA），初始化为2；它不是Wi-Fi中随机抽取退避整数的窗口上限。$BE$决定随机退避区间，$NB$记录本轮接入过程中因信道忙而进行的重试。

每次进入内层，初始化为
\begin{equation}
 NB=0,\qquad CW=2,\qquad BE=BE_{\min}.
\end{equation}
这里$BE_{\min}$为最小退避指数。时隙化系统先定位下一个退避周期（Backoff Period）边界。此处“退避周期”是一个基本时间单位，记其长度为$T_{\mathrm{BP}}$；它与一次完整随机退避持续时间不是同一概念。

\subsubsection{先随机等待，再判断信道}

节点等概率抽取整数
\begin{equation}
 U\sim\mathrm{Unif}\{0,1,\ldots,2^{BE}-1\},
 \qquad T_{\mathrm{bo}}=U T_{\mathrm{BP}}.
\end{equation}
例如$BE=3$时，可抽到0至7，共8个等可能值；若抽到4，就等待4个基本退避周期。取0只表示随机等待为零，不代表绕过后面的CCA直接发送。

本算法在随机等待过程中不采用Wi-Fi式的“忙则冻结、闲则递减”规则。节点先经过所抽出的退避时间，再在规定边界执行CCA。若一次CCA空闲，则$CW$减1；$CW$尚未到0时，继续下一次规定的CCA。只有两次所需的空闲评估均通过，内层才返回发送许可。

若CCA发现信道忙，无论此前已经通过了多少次空闲评估，都重新设置
\begin{equation}
 CW\leftarrow2,\qquad NB\leftarrow NB+1,\qquad
 BE\leftarrow\min(BE+1,BE_{\max}).
\end{equation}
$CW$恢复为2，避免把一次较早的空闲评估与后来跨越忙时段的评估拼接成发送许可；$NB$增加，记录本轮又失败了一次；$BE$增加，但不超过上限，使下一次随机等待的候选范围更宽。若仍有接入重试预算，就重新抽取退避；预算耗尽，则返回接入失败。

\subsubsection{一个包含忙信道的变量演化例子}

设$BE_{\min}=3$，且退避指数上限至少为4，开始时$(NB,CW,BE)=(0,2,3)$。第一次抽到$U=2$，经过两个基本周期后进行CCA。如果该次评估发现信道忙，变量变为$(1,2,4)$。第二次抽样范围于是扩大为0至15；假设抽到$U=3$，节点再等待三个周期。

若随后第一次CCA空闲，$CW$由2变为1；下一次CCA也空闲，$CW$由1变为0，内层返回允许发送。若第二次CCA反而忙，则不能发送，必须把$CW$恢复到2，并再次更新$NB$与$BE$。这一例子只是变量规则的演示，不是对任何节点将抽到哪些随机数的预测。

\begin{figure}[htbp]
\centering
\begin{tikzpicture}[every node/.style={font=\small},x=0.95cm]
\draw[->] (0,0)--(11.3,0) node[right]{时间};
\foreach \x in {0,...,10} {\draw (\x,-0.12)--(\x,0.12);}
\draw (0,0.5) rectangle (2,1.2); \node at (1,0.85) {随机等待2};
\draw (2,0.5) rectangle (3,1.2); \node at (2.5,1.55) {CCA忙};
\draw (3,0.5) rectangle (6,1.2); \node at (4.5,0.85) {新随机等待3};
\draw (6,0.5) rectangle (7,1.2); \node at (6.5,1.55) {CCA闲};
\draw (7,0.5) rectangle (8,1.2); \node at (7.5,2.05) {CCA闲};
\draw (8,0.5) rectangle (11,1.2); \node at (9.5,0.85) {数据发送};
\node at (2.5,-0.6) {$NB=1,\ BE=4$};
\node at (6.5,-0.6) {$CW=1$}; \node at (8,-1.1) {$CW=0$};
\end{tikzpicture}
\caption{随机退避与CCA交替的原创逻辑时序。方框表达过程顺序，CCA方框宽度不作为实际持续时间参数。}
\label{fig:p4-cca}
\end{figure}

\subsection{两类重试预算与半双工时序}

内层的接入重试预算与外层的发送重试预算是两套独立约束。节点在内层遇到忙信道，会增加$NB$，却没有因此增加一次已经发出的数据帧。节点真正发出数据而未收到ACK，才改变外层的发送失败状态。只要外层还允许再次尝试，就重新进入内层，并重新把$NB$初始化为0；新一轮接入重新拥有其规定的退避尝试预算。

因此，一个待发帧可能经历“若干次CCA忙---获准发送---没有ACK---重新进行若干次CCA”这样的过程。分析时必须分别记录“接入了几次”和“真正发送了几次”。最大重试参数还应明确采用“初次尝试之外的重试次数”还是“包含初次尝试的总次数”，不能把两种计数约定直接混用。

\subsubsection{DATA、SIFS、ACK及回复后的恢复}

发送获准之后，数据接收端正确收到帧，经过短帧间间隔（Short Interframe Space, SIFS）再回送ACK。单收发器需要由接收状态转换为发送状态，不能瞬时同时承担两种工作。ACK完成之后，接收端还需从发送切回接收，才能处理下一帧。课程的简化交换时序把两个转换保护阶段都表示为SIFS：
\begin{equation}
 \mathrm{DATA}\;\longrightarrow\;\mathrm{SIFS}\;
 \longrightarrow\;\mathrm{ACK}\;\longrightarrow\;
 \mathrm{SIFS}\;\longrightarrow\;\text{下一次可接收阶段}.
\end{equation}
其核心含义是“没有发射能量”与“接收机已准备好接收”不完全等价。若另一节点恰在回复后的切换阶段开始发送，即使侦听到的无线能量很低，接收端也可能来不及接收新帧开头。

如果只核算这一个简化交换中的占用与保护时间，可以写为
\begin{equation}
 T_{\mathrm{exchange}}=T_{\mathrm{DATA}}+
 2T_{\mathrm{SIFS}}+T_{\mathrm{ACK}}.
\end{equation}
完整交付时间还应加入随机退避、CCA、对齐等待、排队及重传，不能把上式当作从上层产生数据到成功交付的全部时延。实际实现中的帧间隔、收发转换和CCA时序须按对应PHY与标准版本配置；该式服务于本课程的简化半双工分析。

\subsection{无信标模式与非时隙CSMA/CA}

无信标模式（Non-beacon-enabled Mode）不依靠周期信标组织上述同步超帧，采用非时隙CSMA/CA（Unslotted CSMA/CA）。外层仍然是接入许可、数据发送、ACK与失败后的重试；同样存在“没有取得信道”与“已经发送但未获确认”两类失败。

内层的主要简化是不必等候全网共享的退避边界，也不使用时隙方案中初始化为2的$CW$。节点初始化$NB$和$BE$，抽取随机退避并等待；随机等待结束后执行CCA。若评估空闲，返回允许发送；若评估忙，则更新$NB$和$BE$，在仍有预算时重新退避，否则返回接入失败。

“非时隙”不表示没有定时单位，也不表示可以忽略接收端切换时间。退避仍以规定的时间尺度形成随机等待，数据与ACK交换仍需满足相应间隔。与时隙方案一样，退避计数归零本身只是进入侦听步骤，是否发送仍取决于CCA结果。

\subsubsection{非时隙接入的补充数值例：从忙信道到重传成功}

非时隙接入的计算需要同时维护绝对时刻、内层状态与外层状态。以下为依据本节算法构造的补充例题，参数专为完整演示而设，并非原课堂数值题，也不代表标准默认参数。设一帧在$t=0$进入队列，采用
\begin{align}
 T_{\mathrm{BP}}&=100\,\mu\mathrm{s}, &
 T_{\mathrm{CCA}}&=20\,\mu\mathrm{s},\\
 T_{\mathrm{DATA}}&=200\,\mu\mathrm{s}, &
 T_{\mathrm{ACK}}&=40\,\mu\mathrm{s},\\
 BE_{\min}&=2, & BE_{\max}&=3.
\end{align}
取课堂简化模型的$T_{\mathrm{SIFS}}=192\,\mu\mathrm{s}$，分别保护DATA后与ACK后的收发转换。ACK超时从数据发送结束时开始计时，期限取$500\,\mu\mathrm{s}$。该期限长于正常回复所需的$192+40=232\,\mu\mathrm{s}$。

为消除“最大次数”不同命名带来的歧义，本例直接定义阈值：内层每次CCA忙后使$NB$加1，若$NB<3$则重试，达到3则接入失败；外层每次发送后ACK超时使$NF$加1，若$NF<2$则重试，达到2则本帧失败。初始$NB=NF=0$。这一判定方式按课件流程图的严格小于号组织，不据此推定其他标准版本同名参数的含义。

已知发送端第一次CCA判忙，之后列出的CCA均判空闲；第一次数据没有被目的端成功接收，因此不回ACK，第二次数据正确接收。第一次失败原因在本例中作为给定条件，不从“发送端CCA空闲”反推接收端一定没有干扰。

第一步，初始化$NB=0,BE=2$，从$\{0,1,2,3\}$抽到$U=2$。随机等待为$200\,\mu\mathrm{s}$，故在区间$[200,220)$执行CCA。即使退避期间外部信号已经出现，随机等待也继续流逝，不能采用Wi-Fi冻结规则延长这200微秒。到$t=220$得到忙结果，更新为$NB=1,BE=3$。因为$1<3$，继续内层；此时尚未发送数据，故$NF$仍为0。

第二步，从扩大后的$\{0,\ldots,7\}$抽到$U=3$，等待300微秒，故下一次CCA位于$[520,540)$。CCA空闲，数据在$[540,740)$发送。数据没有被正确接收，节点等待至
\begin{equation}
 t_{\mathrm{timeout}}=740+500=1240\,\mu\mathrm{s}
\end{equation}
仍未获得ACK，于是$NF=1$。由于$1<2$，允许再次发送；这时要重新进入内层，设置$NB=0,BE=2$。第一次内层积累的$NB=1$与$BE=3$不能延续到这次重新初始化之后。

第三步，新一轮抽到$U=1$，等待$100\,\mu\mathrm{s}$，在$[1340,1360)$执行CCA。空闲评估通过后，第二次数据在$[1360,1560)$发送。目的端正确收到数据，先切换$192\,\mu\mathrm{s}$，因此ACK开始于
\begin{equation}
 t_{\mathrm{ACK,start}}=1560+192=1752\,\mu\mathrm{s},
\end{equation}
并于$1752+40=1792\,\mu\mathrm{s}$结束。ACK及时到达，外层以成功结束；发送端从$t=0$到收到完整确认共经过1792微秒。

目的端的状态还需要继续向后画：在ACK结束后，它经过192微秒切回接收，至
\begin{equation}
 t_{\mathrm{ready}}=1792+192=1984\,\mu\mathrm{s}
\end{equation}
才按本例模型准备好接收下一帧。因此，“本帧已获确认”的1792微秒与“汇聚节点可接收下一帧”的1984微秒是两个不同事件。

\begin{figure}[htbp]
\centering
\begin{tikzpicture}[every node/.style={font=\small}]
\draw[->] (0,0)--(11.4,0) node[right]{$t$};
\draw (0,0.4) rectangle (2.1,1.2); \node[align=center] at (1.05,0.8) {退避、忙CCA\\再退避};
\draw (2.1,0.4) rectangle (4,1.2); \node[align=center] at (3.05,0.8) {CCA闲\\DATA失败};
\draw (4,0.4) rectangle (6.1,1.2); \node[align=center] at (5.05,0.8) {等待ACK\\超时};
\draw (6.1,0.4) rectangle (8.5,1.2); \node[align=center] at (7.3,0.8) {内层重置\\再接入与DATA};
\draw (8.5,0.4) rectangle (10,1.2); \node[align=center] at (9.25,0.8) {SIFS\\ACK};
\draw[dashed] (10,0.4) rectangle (11.2,1.2); \node at (10.6,0.8) {SIFS};
\foreach \x/\v in {0/0,4/740,6.1/1240,8.5/1560,10/1792,11.2/1984}{\draw (\x,-0.1)--(\x,0.2); \node[rotate=45,anchor=east] at (\x,-0.2) {$\v$};}
\node at (5.6,-1.25) {时间标签单位：微秒；图形宽度不按比例};
\end{tikzpicture}
\caption{补充非时隙例题的原创过程图。接入忙、数据失败和接收端恢复被分开表示。}
\label{fig:p4-unslotted-example}
\end{figure}

这个结果还可以按时间分量核对。第一次发送前耗时540微秒，第一次数据及超时耗时$200+500=700$微秒；第二轮退避、CCA、DATA、SIFS与ACK耗时
\begin{equation}
 100+20+200+192+40=552\,\mu\mathrm{s}.
\end{equation}
合计$540+700+552=1792\,\mu\mathrm{s}$，与绝对时间轴一致。成功之后的192微秒恢复时间没有计入发送端获得确认的时延，却必须计入目的端再次可接收的时刻。

\subsubsection{同一规则下的两种丢包分支}

如果内层连续三次CCA忙，在本例阈值下，$NB$依次为1、2、3；第三次忙后退出为接入失败。假设初始退避抽样依次为$U=2,3,0$，则三次CCA分别位于$[200,220)$、$[520,540)$、$[540,560)$，在$t=560$结束为接入失败。第三次抽到0不免除CCA；$BE$已达到上限3，不再继续增大。这个分支中数据发送次数为零，$NF$仍为0，外层却仍然结束为失败。

另一种情况是每次内层都成功取得信道，但连续两次实际发送后都ACK超时。第一次超时后$NF=1$，允许新一轮接入；第二次超时后$NF=2$，达到本例阈值，停止重传。此时最后一轮的$NB$即使为0，也不能抵消外层已经耗尽的发送预算。

若另一节点在成功例子的$t=1800\,\mu\mathrm{s}$开始发送给同一个汇聚节点，按本例简化模型它落在ACK后的恢复区间$[1792,1984)$中。接收端尚未恢复，因此新帧可能丢失。这个分支说明，数值题必须同时绘制发送端的接入状态和接收端的收发状态，不能只看一条代表无线能量的忙闲线。

\subsection{返回DCF：统一原则与不同实现}

\subsubsection{CSMA/CA不是某一套固定流程图}

载波侦听多路访问与碰撞避免（CSMA/CA）是一类接入思想：发送前评估介质状态，在不能安全接入时推迟，并通过随机化等机制减少同时发送。DCF是Wi-Fi中这一思想的一种具体实现，IEEE 802.15.4又给出另一套具体规则。说“DCF加退避”容易造成误解，因为退避本来就是DCF的一部分。

Wi-Fi的退避通常在允许递减的空闲时隙下降，信道忙时冻结，待满足恢复条件后继续剩余值。这里讨论的802.15.4基本流程则先经过随机退避，再执行CCA，遇忙重新抽取；两者不能只因都出现“backoff”就照搬计数规则。Wi-Fi的$CW$描述抽样窗口，时隙802.15.4的$CW$描述剩余空闲评估次数，也不能互相代入公式。

\subsubsection{NAV、冻结和推迟分别表达什么}

NAV是一种虚拟载波侦听（Virtual Carrier Sensing）状态，根据正确收到的Duration信息保留预期交换时间。物理载波侦听（Physical Carrier Sensing）则观察无线介质是否存在被判定为忙的活动。对DCF的接入判断，可以把两者关系概括为
\begin{equation}
 \mathrm{Busy}=\mathrm{PHYBusy}\ \lor\ (NAV>0).
\end{equation}
因此，NAV为零不保证可以递减退避；只要物理层判忙，仍然需要冻结。反之，即使当下没有检测到发射能量，NAV尚未到期，也不能把这段时间当成自由竞争机会。

冻结（Freeze）指保存退避计数器的当前值，暂不递减；推迟（Defer）是发送没有立即发生的行为结果，不是另一个与NAV并列的硬件计数器。设置NAV可能导致冻结，但冻结并不要求必定先设置NAV。例如隐藏站点听到AP发送ACK时，即便没有建立新的正NAV，也能因物理忙而冻结退避。

恢复退避一般仍需满足相应帧间空闲条件，不能在忙信号刚结束的瞬间就任意恢复时隙递减。是否需要退避，又由待发帧遇忙、失败后重传、成功后的后退避等状态决定。单独出现一段等待时间，不足以判断节点必然在其后新抽一个退避数。

\subsubsection{EIFS保护未知交换，RTS门限依据帧长}

EIFS适用于不能正确解码已接收活动、因而无法可靠获取交换持续时间的保护情形。两个站点同时发送，旁观站点可能只看到物理忙而无法读取首部；忙结束后，它用EIFS保护潜在的后续回复。如果等待期间正确接收到新的帧，节点可依据新的有效信息回到正常的接入规则。

不同包长的碰撞尤其能显示发送端之间的差异：短帧发送端结束后只听到长帧尾部，可能无法正确解码，从而采用EIFS；长帧发送端则主要依据自身ACK超时判断失败。EIFS在适用时替代通常的帧间保护等待，不应机械地再叠加一个完整DIFS；是否还需执行退避，则取决于节点已有的DCF状态。

请求发送/清除发送（Request to Send/Clear to Send, RTS/CTS）通过短控制交换帮助预留介质，能够减轻某些隐藏节点造成的长帧碰撞代价。在基于门限的选择中，是否使用RTS/CTS由待发帧长度与RTS门限比较决定，不能简单写成“发现隐藏节点就一定使用RTS”。保护机制的收益还需与控制报文开销比较。这一结论把本部分的传感器接入过程与前面的Wi-Fi机制重新连成一条主线：相同的共享介质问题，可以因设备能力、能量目标和业务规模不同而采用不同具体规则。

\subsection*{本部分资料说明}
\addcontentsline{toc}{subsection}{本部分资料说明}
本部分按10月14日上午、10月14日后半段、10月15日两段课程和10月16日上午的授课顺序组织；其中10月15日开头的隐藏节点解答与前一天题设合并，以保持推导连续。图均为根据概念与过程重新设计的教学示意，不作为原始课堂板书复原。隐藏节点例题采用已明确给出的归一化时隙关系；具体标准的边界参数与版本差异不由课堂简化模型替代。电网监测仅用于说明感知与控制需求，不据课堂案例判定某次实际停电的原因。
